传统媒体与网络平台的交替发展，使艺人的职业路径与公众形象呈现出跨媒介变化。本文以佐佐木希为主题，构造一个覆盖2002年至2025年的虚构纵向案例，围绕职业阶段划分、媒介形象编码与公众情感变化建立分析流程。将职业轨迹划分为时尚模特、影视转型、多领域拓展、职业稳定与新媒体重塑五个阶段，结合人物图片讨论视觉表达，再采用加权情感得分与编码一致性指标，对四个时间窗口的模拟评论进行分析。

模拟分析得到的四个情感得分依次为0.25、0.40、0.10和0.55，反映职业重塑阶段内部的评价波动。分析表明，长期职业定位与短期互动反馈应在不同时间尺度上解释，图片中的视觉风格也需要结合媒介场景理解。本文的时间线、事件、样本与数值均为教学设定，不构成对真实人物经历或公众评价的事实结论。后续研究可通过跨平台抽样、独立编码与多案例比较检验分析框架的适用性。
